\chapter{Operações Lógicas  II}

As operações lógicas obedecem às regras do cálculo proposicional e são fundamentais para a construção e análise de proposições mais complexas.

Entre as principais operações lógicas, destacam-se:
\begin{itemize}
    \item Negação;
    \item Conjunção;
    \item Disjunção;
    \item Disjunção exclusiva;
    \item Condicional;
    \item Bicondicional.
\end{itemize}

Neste capítulo, aprofundaremos o estudo das operações de disjunção exclusiva, condicional e bicondicional.

\section{Disjunção Exclusiva}

\subsection{Definição}

Chama-se de \textbf{disjunção exclusiva} de duas proposições $p$ e $q$ a proposição representada por ``ou $p$ ou $q$'', no sentido de que \textbf{exatamente uma} das proposições é verdadeira.

\subsection{Notação}

A disjunção exclusiva é representada por:
\[
p \ \underline{\vee} \ q
\]

\subsection{Tabela-Verdade}

\begin{center}
\begin{tabular}{|c|c|c|}
\hline
$p$ & $q$ & $p \ \underline{\vee} \ q$ \\
\hline
V & V & F \\
V & F & V \\
F & V & V \\
F & F & F \\
\hline
\end{tabular}
\end{center}

\subsection{Exemplo}

Considere as proposições:
\begin{itemize}
    \item $p$: Rivaldo foi revelado no Santa Cruz. (V)
    \item $q$: Neymar é melhor do que Cristiano Ronaldo. (F)
\end{itemize}

A proposição composta:
\[
p \ \underline{\vee} \ q
\]
pode ser lida como:

\begin{quote}
Ou Rivaldo foi revelado no Santa Cruz ou Neymar é melhor do que Cristiano Ronaldo.
\end{quote}

Como $p$ é verdadeiro e $q$ é falso, temos:
\[
V(p \ \underline{\vee} \ q) = V
\]

\section{Condicional}

\subsection{Definição}

Chama-se de \textbf{proposição condicional} (ou simplesmente \textbf{condicional}) a proposição representada por:

\begin{quote}
Se $p$ então $q$
\end{quote}

\subsection{Notação}

\[
p \rightarrow q
\]

\subsection{Tabela-Verdade}

\begin{center}
\begin{tabular}{|c|c|c|}
\hline
$p$ & $q$ & $p \rightarrow q$ \\
\hline
V & V & V \\
V & F & F \\
F & V & V \\
F & F & V \\
\hline
\end{tabular}
\end{center}

\subsection{Exemplo}

Considere:
\begin{itemize}
    \item $p$: Rivaldo foi revelado no Santa Cruz. (V)
    \item $q$: Neymar é melhor do que Cristiano Ronaldo. (F)
\end{itemize}

A proposição:
\[
p \rightarrow q
\]
é lida como:

\begin{quote}
Se Rivaldo foi revelado no Santa Cruz, então Neymar é melhor do que Cristiano Ronaldo.
\end{quote}

Nesse caso:
\[
V(p \rightarrow q) = F
\]

\subsection{Outras Formas de Leitura}

A proposição $p \rightarrow q$ também pode ser interpretada como:
\begin{itemize}
    \item $p$ é condição suficiente para $q$;
    \item $q$ é condição necessária para $p$.
\end{itemize}

\subsection{Observação Importante}

A proposição $p \rightarrow q$ \textbf{não afirma} que $q$ pode ser deduzido de $p$. Trata-se apenas de uma relação lógica entre os valores de verdade das proposições.

\section{Bicondicional}

\subsection{Definição}

Chama-se de \textbf{proposição bicondicional} (ou simplesmente \textbf{bicondicional}) a proposição representada por:

\begin{quote}
$p$ se e somente se $q$
\end{quote}

\subsection{Notação}

\[
p \leftrightarrow q
\]

\subsection{Tabela-Verdade}

\begin{center}
\begin{tabular}{|c|c|c|}
\hline
$p$ & $q$ & $p \leftrightarrow q$ \\
\hline
V & V & V \\
V & F & F \\
F & V & F \\
F & F & V \\
\hline
\end{tabular}
\end{center}

\subsection{Exemplo}

Considere:
\begin{itemize}
    \item $p$: Rivaldo foi revelado no Santa Cruz. (V)
    \item $q$: Neymar é melhor do que Cristiano Ronaldo. (F)
\end{itemize}

A proposição:
\[
p \leftrightarrow q
\]
é lida como:

\begin{quote}
Rivaldo foi revelado no Santa Cruz se e somente se Neymar é melhor do que Cristiano Ronaldo.
\end{quote}

Como os valores de $p$ e $q$ são diferentes, temos:
\[
V(p \leftrightarrow q) = F
\]

\subsection{Outras Formas de Leitura}

A proposição bicondicional pode ser interpretada como:
\begin{itemize}
    \item $p$ é condição necessária e suficiente para $q$;
    \item $q$ é condição necessária e suficiente para $p$.
\end{itemize}